% ====================================================================
% SECTION 06: THE ARCHITECTURAL TRANSITION
% ====================================================================

\section{The Architectural Transition}
\label{sec:womersley_minimax}

The gauge-invariant network Lagrangian established in \S\ref{sec:gauge} assumes
a scale-invariant wave penalty $\mathcal{C}_{\mathrm{wave}}$. This is exact only
in the limit $\mathrm{Wo}\to\infty$, where fluid inertia dominates and the
impedance of a vessel depends solely on its cross-sectional area. In the general
case, the reflection at each junction is governed by the mismatch of the complex
Womersley admittances $Y(\mathrm{Wo}) \propto R^2 \sqrt{1-F_{10}(\mathrm{Wo})}$.

\begin{theorem}[Analytic Admittance Matching]
The universal vascular branching exponent $\alpha$ is bounded by two exact
fluid-dynamic attractors representing the limits of zero reflection ($\Gamma =
0$):
\begin{enumerate}
    \item
\textbf{Inertial Attractor ($\mathrm{Wo} \to \infty$):} $F_{10} \to 0$, thus $Y
\propto R^2$. Impedance matching requires $R_p^2 = \sum R_{c,i}^2$, yielding the
area-preserving law $\alpha = 2$.
    \item
\textbf{Viscous Attractor ($\mathrm{Wo} \to 0$):} $1-F_{10} \propto
\mathrm{Wo}^2$, thus $\sqrt{1-F_{10}} \propto \mathrm{Wo}$. Since $\mathrm{Wo}
\propto R$, the admittance scales as $Y \propto R^3$. Impedance matching
requires $R_p^3 = \sum R_{c,i}^3$, yielding Murray's Law $\alpha = 3$.
\end{enumerate}
\end{theorem}

The allometric transition is thus the trajectory of the network-level minimax as
the characteristic Womersley number scales with body mass $M$. To ensure a
\chA{physically constrained derivation}, the wave penalty
$\mathcal{C}_{\mathrm{wave}}$ is defined as the spectral sum of reflected power
across the first five harmonics $\{n\omega_H\}$ of the cardiac cycle, weighted
by the longitudinal $k$-dispersion attenuation $e^{-2 \kappa L}$ (see
Supplemental Material, Section S2):
\begin{equation}
    \mathcal{C}_{\mathrm{wave}} = \sum_{n=1}^5 H_n \sum_{g=0}^{G-1} \left( \prod_{j=0}^{g} \eta_{\mathrm{att}, j}^{(n)} \right) |\Gamma_g^{(n)}|^2,
\end{equation}
where $H_n$ are the fixed power weights of the mammalian pulse spectrum. This
formulation requires no ad-hoc damping functions or scaling constants.
Sensitivity analysis (see Supplemental Material, Section S4) demonstrates that
$\alpha^*$ is robust to variations in the spectral envelope, as the fundamental
harmonic $H_1$ dominates the energetic partition. The 'viscous shielding' of the
microcirculation emerges naturally from the exponential decay of oscillatory
energy in the low-$\mathrm{Wo}$ regime. The critical transition at $M^* \approx
\VarMStar\,$g is the physical threshold where viscous dissipation cedes
dominance to the incommensurability principle. This threshold is governed by the
scaling of the characteristic Womersley number:
\begin{equation}
    \mathrm{Wo}_g \approx \mathrm{Wo}_{\mathrm{ref}} \left( \frac{r_g}{r_{\mathrm{ref}}} \right) \sqrt{\frac{f}{f_{\mathrm{ref}}} \frac{\nu_{\mathrm{ref}}}{\nu}},
    \label{eq:wo_scaling}
\end{equation}
where $\mathrm{Wo}_{\mathrm{ref}}$ is the value at a known physiological scale
(e.g., human adult aorta).

\begin{theorem}[The Cross-Class Allometric Transition]
\label{thm:cross_transition}
While the minimax attractor $\alpha^*$ is structurally invariant within a fixed
allometric class (Theorem~\ref{thm:arch}), the transition between distinct
regimes (viscous vs. wave-dominated) is governed by the crossing of a critical
Womersley interval $[\mathrm{Wo}_c^{\mathrm{fluid}},
\mathrm{Wo}_c^{\mathrm{wave}}] = [\sqrt{3}, 3/\sqrt{2}]$. This dual-threshold
transition induces a sharp allometric shift in branching geometry at a body mass
$M^* \approx \VarMStar\,$g.
\end{theorem}

\begin{theorem}[Emergent Allometric Transition Limits]
\label{thm:emergent_transition}
For a given body mass $M$, let $\alpha^*(M)$ be the unique minimax saddle point
of $\mathcal{L}_{\mathrm{net}}(\alpha,\eta; M)$. Then the trajectory of the
attractor obeys the following limits:
\begin{enumerate}
    \item
As $M\to 0$, $\alpha^*(M) \to \alpha_t \approx \chA{\VarAlphaT}$ (viscous
dominance).
    \item
As $M\to\infty$, $\alpha^*(M) \to \alpha_w \approx 2.0$ (inertial dominance).
    \item
There exists a unique mass $M^* \approx \VarMStar\,$g where the slope
$|d\alpha^*/d(\ln M)|$ is maximal, defining the sigmoidal inflection point of
the allometric transition.
\end{enumerate}
\end{theorem}

\begin{proof}
The limits follow from the asymptotic behavior of the admittance
$Y(\mathrm{Wo})$. For $M\to 0$, $\mathrm{Wo}\to 0$ throughout the tree, the
reflection coefficient tends to zero, removing the wave penalty and allowing the
transport cost to pull $\alpha^*$ toward its static minimum at $\alpha_t \approx
\chA{\VarAlphaT}$. For $M\to\infty$, the inertial geometric matching condition
$\alpha_w \approx 2.0$ dominates. The sigmoidal transition in
Fig.~\ref{fig:transition} emerges directly from this numerical competition,
smoothly connecting the two topological attractors.
\end{proof}

\begin{corollary}[Fourth-Power Scaling of Transition Mass]
\label{cor:fourth_power}
The critical mass $M^*$ at which the allometric transition occurs scales as the
fourth power of the critical Womersley number:
\begin{equation}
    M^* \propto \mathrm{Wo}_c^4.
    \label{eq:fourth_power}
\end{equation}
For networks embedded in different dimensions, this yields a topologically
grounded prediction:
\begin{equation}
    \frac{M^*_{d=2}}{M^*_{d=3}} = \left(\frac{\mathrm{Wo}_c(d{=}2)}{\mathrm{Wo}_c(d{=}3)}\right)^4
    = \left(\frac{\sqrt{6}}{\sqrt{3}}\right)^4 = 4.
    \label{eq:dimension_ratio}
\end{equation}
\end{corollary}

\begin{proof}
The transition mass is defined by the condition that the characteristic
Womersley number at a reference generation (e.g., the aorta) reaches the
critical fluid threshold $\mathrm{Wo}_c^{\mathrm{fluid}}$. From the allometric
scaling of heart rate $f \propto M^{-1/4}$ and aortic radius $r_0 \propto
M^{3/8}$ (derived from Kleiber's Law via flow continuity), the dimensional form
of the Womersley number gives:
\begin{equation}
    \mathrm{Wo}_0 = r_0 \sqrt{\frac{2\pi f}{\nu}}
    \propto M^{3/8} \cdot M^{-1/8} = M^{1/4}.
\end{equation}
Setting $\mathrm{Wo}_0(M^*) = \mathrm{Wo}_c$ yields the scaling relation $M^*
\propto \mathrm{Wo}_c^4$. However, this dimensional estimate (which would
predict $M^* \sim 5$--$10\,$g for typical mammalian parameters) represents only
an order-of-magnitude threshold. The precise value $M^* \approx \VarMStar\,$g is
determined numerically as the inflection point of the sigmoidal transition
$\alpha^*(M)$, i.e., the body mass at which $|d\alpha^*/d(\ln M)|$ is maximal
(Theorem~\ref{thm:emergent_transition}, item~3). This sharper definition
captures the smooth, progressive nature of the allometric crossover, rather than
a discontinuous jump at the Womersley threshold.

The dimensional ratio \eqref{eq:dimension_ratio} follows immediately from the
kinematic matching criterion applied to the bulk fluid
$\mathrm{Wo}_c^{\mathrm{fluid}}(d) = \sqrt{6/(d-1)}$ derived in
\S\ref{sec:isotropic_stress}. This prediction is experimentally testable by
comparing the allometric transitions in planar networks (e.g., retinal
vasculature) versus three-dimensional networks (e.g., coronary circulation).
\end{proof}

\subsection{Kinematic Matching Criterion at the Bifurcation}
\label{sec:isotropic_stress}

The architectural transition at $M^*$ marks the critical threshold where the
vascular tree can no longer ignore the complex phase of the fluidic impedance.
To determine this critical Womersley number $\mathrm{Wo}_c$, we must address a
fundamental epistemological boundary in biological transport modeling.

The 1D Navier-Stokes equations governing pulsatile flow within a single vascular
segment yield the complex Womersley admittance $Y(\mathrm{Wo})$. However, these
local equations of motion are \emph{topologically blind}: they contain no
information regarding the dimensionality $d$ of the space the network is
required to fill. Consequently, it is mathematically impossible to derive a
global geometric threshold strictly from the local fluid dynamics. For the 1D
fluidic transport to stably support a space-filling hierarchy, the local
hydrodynamics must geometrically couple with the $d$-dimensional structural
embedding.

Rather than relying on phenomenological fitting, we resolve this
mechanical-topological decoupling by applying the \textbf{Kinematic Matching
Criterion} (Theorem~\ref{thm:kinematic_matching}). This criterion establishes a
\chA{physical matching condition}---combining Euclidean geometry, classical
kinematics, and evanescent mode absorption at the junction---demonstrating that
the pulsatile network undergoes a critical transition when the viscous
absorption capacity of the Womersley flow (quantified by the inverse quality
factor $\mathcal{Q}^{-1}$) exactly balances the geometric scattering imposed by
the space-filling topology. For the network to maintain this equilibrium, the
inverse quality factor must satisfy:
\begin{equation}
    \mathcal{Q}^{-1} = d - 1.
    \label{eq:isotropic_condition}
\end{equation}

This criterion acts as the essential boundary condition bridging 1D transport
and $d$-dimensional topology. For mammalian arterial beds embedded in 3D tissue
($d=3$), this requires $\mathcal{Q}^{-1} = 2$. By combining the Navier-Stokes
expression for viscous absorption ($\mathcal{Q}^{-1} = 6/\mathrm{Wo}^2$) with
the Euclidean space-filling geometry, the criterion provides the \chA{physical
foundation} for the allometric transition.

However, pulsatile transport involves two mathematically distinct physical
quantities: the \emph{local fluid mass} (governed by the longitudinal admittance
$Y_L$) which undergoes viscous dissipation, and the \emph{propagating pressure
wave} (governed by the characteristic admittance $Y_c$) which \chA{governs}
network reflections. Because \chA{$Y_c \propto \sqrt{i\omega\rho\,Y_L}$,
equivalently $Y_c\propto\sqrt{1-F_{10}}$}, the wave phase is algebraically
halved with respect to the bulk fluid, inducing a profound \textbf{Phase
Decoupling}. Applying the kinematic matching criterion ($\mathcal{Q}^{-1} =
d-1$) to both quantities yields two distinct thresholds (derived in Supplemental
Section S1). \chA{We emphasise that the isotropic-embedding criterion
$\mathcal{Q}^{-1}=d-1$ is derived in S1 for the longitudinal (bulk-fluid) mode;
its application to the characteristic admittance $Y_c$ is a
\emph{postulate}---we extend the same kinematic matching to the propagating wave
by analogy, with the algebraic phase-halving
\chA{$Y_c\propto\sqrt{i\omega\rho\,Y_L}$, equivalently
$Y_c\propto\sqrt{1-F_{10}}$} as its physical basis. We adopt this as a working
hypothesis rather than a derived result; the qualitative conclusion (two
incommensurable thresholds) follows from the phase-halving alone and is robust
to the precise value of the wave threshold.}

\chRB{The intuition behind the two thresholds is that a bifurcation must cope
with \emph{two} things at once: the bulk fluid it carries, and the pressure wave
riding on that fluid. As an animal grows and its Womersley number rises, the
bulk fluid is the first to ``switch on''---to stop dissipating its junction kick
viscously and begin reflecting it---at $\mathrm{Wo}_c^{\mathrm{fluid}}$. The
pressure wave, because it responds to the \emph{square root} of the fluid
admittance, lags behind and switches on only at the larger
$\mathrm{Wo}_c^{\mathrm{wave}}$. The two can therefore never be perfectly
matched at the same body size: whatever the animal's scale, one channel is
always on the ``wrong'' side of its own threshold. That permanent, unavoidable
mismatch is precisely what we mean by incommensurability, and it is why no
single branching exponent can zero out both costs at once.}

\begin{enumerate}
    \item
\textbf{The Fluid Threshold ($Y_L$):} The local bulk fluid achieves
space-filling resonance when its dissipation ratio matches the topology:
    \begin{equation}
        \frac{6}{\mathrm{Wo}^2} = d - 1 \quad \Rightarrow \quad \mathrm{Wo}_c^{\mathrm{fluid}} = \sqrt{\frac{6}{d-1}}.
    \end{equation}
For $d=3$, this yields $\mathrm{Wo}_c^{\mathrm{fluid}} = \sqrt{3} \approx
1.732$.
    \item
\textbf{The Wave Threshold ($Y_c$):} The propagating wave, subjected to the
square root of the local admittance, achieves kinematic matching at a higher
Womersley number:
    \begin{equation}
        \frac{\sqrt{36+\mathrm{Wo}^4}+\mathrm{Wo}^2}{6} = d - 1 \quad \Rightarrow \quad \mathrm{Wo}_c^{\mathrm{wave}} = \sqrt{\frac{3d(d-2)}{d-1}}.
    \end{equation}
For $d=3$, this yields $\mathrm{Wo}_c^{\mathrm{wave}} = \frac{3}{\sqrt{2}}
\approx 2.121$.
\end{enumerate}

The fact that these two thresholds do not coincide constitutes the exact
mathematical formulation of the \textbf{Incommensurability Principle}: the fluid
and the wave can never achieve equipartition simultaneously. The ratio of their
squared thresholds, $\frac{d(d-2)}{2}$, is never equal to 1 for any integer $d$.
Remarkably, Womersley's historical empirical transition threshold
($\mathrm{Wo}_c \approx 2.0$, empirically utilized in Paper~III) falls precisely
within the theoretical gap $[\sqrt{3}, 3/\sqrt{2}]$ defined by this decoupling,
\chA{connecting the phenomenological observations with the physical} theory.

\chA{For clarity, the \emph{quantitative} predictions of this work---the
transition mass $M^*$, the $\alpha^*(M)$ curve, and the per-species
exponents---enter the solver only through the fluid threshold
$\mathrm{Wo}_c^{\mathrm{fluid}}=\sqrt{3}$, which sets the viscous-to-inertial
crossover used to compute $\alpha^*(M)$. The wave threshold
$\mathrm{Wo}_c^{\mathrm{wave}}=3/\sqrt{2}$ fixes the \emph{upper} edge of the
incommensurability gap and the qualitative regime structure (and the $d=2$
retinal collapse), but it does not appear in the numerical optimization. The
incommensurability of the two thresholds is thus a structural, regime-defining
statement; the headline numbers depend only on $\sqrt{3}$.}

The predictive power of this Phase Decoupling is \chA{supported} by its global
consequences. First, it defines a \textbf{scaling transition interval} in body
mass $M^* \sim 1\,$g (plausible range $0.5\text{--}1.2\,$g depending on
species-specific heart rate variability), separating viscous-dominated from
wave-influenced mammalian architectures. Second, it provides a definitive
structural explanation for the ``Retinal Paradox'' (Mechanical-Fluidic
Decoupling, \S\ref{sec:retinal}): for networks constrained to a 2D manifold
(e.g., the human retina, $d=2$), the wave threshold collapses exactly to zero:
$\mathrm{Wo}_c^{\mathrm{wave}}(d{=}2) = 0$. This implies that in a 2D topology,
the propagating wave is \emph{permanently} above its threshold
($\mathcal{Q}^{-1}_{Y_c} \ge 1$ for all $\mathrm{Wo} > 0$). \chA{Here the
criterion is applied to the characteristic (transmitted) admittance $Y_c$, for
which a large $\mathcal{Q}^{-1}_{Y_c}$ denotes a permanently impedance-matched,
propagating wave channel---the opposite sense to the longitudinal-mode damping
convention, in which a large $\mathcal{Q}^{-1}$ signals overdamping; the two
senses must not be conflated.} This geometric constraint forces planar retinal
diameters to permanently transition toward the wave-attractor limit ($\alpha
\approx 2.0$), \chA{supporting the interpretation} that the branching geometry
is topologically anchored by the embedding dimension rather than local metabolic
fine-tuning.

\textbf{Falsifiable Milestones.} The theory predicts that organisms below $M^*
\approx \VarMStar\,$g (sub-gram invertebrates, early neonates) should exhibit
branching exponents closer to the Murray attractor ($\alpha \approx 3$), while
all vascular mammals above this threshold should share the wave-influenced
attractor ($\alpha \approx 2.5$--$2.8$). This prediction is testable via
morphometric analysis of vascular casts across developmental stages.

\begin{table}[H]
\centering
\caption{\textbf{Womersley Transition scaling ($\beta$) vs Body Mass.}
Comparison of the theoretical branching parameter $\beta$ predicted by
the emergent minimax framework. The theoretical
model reproduces the architectural transition from viscous-dominated
morphologies ($\beta \approx 0.85$) to the pulsatile attractor
($\beta \approx \VarBetaHumanObs$) \chA{without fitting the target exponent}. The simplified model predicts
a nearly constant branching ratio $\beta \approx \VarBetaHumanObs$ for all mammals; the
observed variation (see, e.g., \cite{kassab1993,huang1996}) requires the
inclusion of taper and asymmetry as in Paper~II.}
\vspace{0.5em}
\begin{tabular*}{\textwidth}{@{\extracolsep{\fill}} l r c c @{}}
\toprule
Species & Mass ($M$) & $\beta$ (Phenomenological) &
$\beta$ (Emergent Minimax) \\
\midrule
Shrew      & $\VarMassShrew\,$g   & \VarBetaShrewObs & \VarBetaShrewWocThree \\
Mouse      & $\VarMassMouseHuo\,$g  & \VarBetaMouseObs & \VarBetaMouseHuoWocThree \\
Rat        & $\VarMassRat\,$g & \VarBetaRatObs & \VarBetaRatWocThree   \\
Guinea Pig & $\VarMassGuineaPig\,$g & \VarBetaGuineaPigObs & \VarBetaGuineaPigWocThree\\
Rabbit     & $\VarMassRabbit\,$g  & \VarBetaRabbitObs & \VarBetaRabbitWocThree\\
Human      & $70\,\mathrm{kg}$  & \VarBetaHumanObs & \VarBetaHumanWocThree \\
Horse      & $500\,\mathrm{kg}$ & \VarBetaHorseObs & \VarBetaHorseWocThree \\
\bottomrule
\end{tabular*}
\label{tab:beta_comparison}
\end{table}

\begin{figure}[H]
\centering
\begin{tikzpicture}
\begin{semilogxaxis}[
    width=\columnwidth,
    height=0.70\columnwidth,
    xlabel={Body Mass $M$ (g)},
    ylabel={Branching Exponent $\alpha^*$},
    xmin=0.1, xmax=10000000,
    ymin=1.8, ymax=3.2,
    grid=both,
    grid style={line width=.1pt, draw=gray!10},
    legend style={at={(0.5,-0.18)}, anchor=north,
                  legend columns=-1, font=\footnotesize},
    ytick={1.8, 2.0, 2.2, 2.4, 2.8, 3.0, 3.2},
    tick label style={font=\small},
    label style={font=\normalsize},
    extra y ticks={2.50, 2.60, 2.70, 2.75, 2.85, 2.95},
    extra y tick style={
        grid=none,
        tick label style={font=\tiny, color=blue!70!black, xshift=-1pt}
    }
]

% Heterogeneity band (symmetric minimax +/- 0.10)
\addplot[name path=upper, draw=none, forget plot]
    table [x=Mass, y expr=\thisrow{Alpha}+0.10]
    {figures/transition_alpha.dat};
\addplot[name path=lower, draw=none, forget plot]
    table [x=Mass, y expr=\thisrow{Alpha}-0.10]
    {figures/transition_alpha.dat};
\addplot[fill=blue!12, forget plot]
    fill between[of=upper and lower];

\addlegendimage{area legend, fill=blue!12, draw=none}
\addlegendentry{Heterogeneity Corridor}

% Symmetric minimax theory curve
\addplot[blue, line width=1.2pt]
    table [x=Mass, y=Alpha] {figures/transition_alpha.dat};
\addlegendentry{Symmetric Minimax}

% Transport (viscous) attractor: alpha_t
\addplot[green!60!black, dashed, domain=0.1:10000000]
    {\VarAlphaTFig};
\addlegendentry{Transport $\alpha_t$}

% Wave (inertial) attractor: alpha_w = 2
\addplot[purple!70, dashed, domain=0.1:10000000]
    {\VarAlphaW};
\addlegendentry{Wave $\alpha_w$}

% Transition mass M* vertical marker
\addplot[gray, densely dotted, line width=0.8pt]
    coordinates {(\VarMStar, 1.8) (\VarMStar, 3.2)};
\node[gray, font=\footnotesize, anchor=center, fill=white, inner sep=2pt, rotate=90]
    at (axis cs:\VarMStar, 2.5)
    {$M^* \approx \VarMStar\,$g};

% Horizontal dashed guides for landing values at M = 10^7 g
\draw[blue!35, densely dashed, line width=0.4pt]
    (axis cs:0.1, 2.50) -- (axis cs:10000000, 2.50);
\draw[blue!35, densely dashed, line width=0.4pt]
    (axis cs:0.1, 2.60) -- (axis cs:10000000, 2.60);
\draw[blue!35, densely dashed, line width=0.4pt]
    (axis cs:0.1, 2.70) -- (axis cs:10000000, 2.70);

\end{semilogxaxis}
\end{tikzpicture}

\caption{\textbf{The Incommensurability Jump: Allometric Transition
    of Branching Geometry.}
    The emergent branching exponent $\alpha^*(M)$ as a function
    of body mass, computed from the symmetric minimax saddle
    point \chA{with no target-exponent fit}. Horizontal dashed lines mark
    the two attractors: the transport limit
    $\alpha_t \approx \VarAlphaTFig$ ($\eta \to 0$) and the
    wave limit $\alpha_w = \VarAlphaW$ ($\eta \to 1$\chA{; this is the
    rigid-wall baseline used to compute the symmetric minimax curve,
    distinct from the elastic-corrected $\alpha_w = \VarAlphaWFig$ shown
    in the conceptual phase diagram, Fig.~\ref{fig:phase_diagram}}).
    The shaded band ($\pm 0.10$) represents the estimated
    range of morphometric heterogeneity corrections
    (bifurcation asymmetry, taper; see \S\ref{sec:gap}).
    The vertical dotted line marks the predicted transition
    mass $M^* \approx \VarMStar\,$g.
    Quantitative per-species validation is provided in
    Table~\ref{tab:beta_comparison}.}
\label{fig:transition}
\end{figure}

\subsection{Meta-Analysis of the Allometric Transition}

The transition observed in Figure~\ref{fig:transition} suggests that mammalian
evolution is governed by a "metacritical" boundary at $M^* \approx
\VarMStar\,$g. This threshold is not merely a statistical centroid but
corresponds to the biological transition from viscous-dominated transport to
wave-dominated pulsatility.

The simplified model places the transition at $M^* \approx \VarMStar\,$g, the
boundary between sub-gram organisms and the smallest mammals. For all mammals
above this threshold, the minimax attractor settles near $\alpha^* \approx
\VarAlphaStarModel$, with only weak residual mass dependence. The full
morphometric analysis of Paper~II, which incorporates generation-by-generation
vessel taper and asymmetry, is required to shift the transition to the 30--60~g
range where rodent weaning occurs; this remains a quantitative refinement for
future work. \chA{We stress that the falsifiable predictions below are anchored
on the Womersley number relative to $\mathrm{Wo}_c = \sqrt{3}$, not on the
absolute transition mass: $M^*$ is a model-dependent landmark---the single-scale
estimate $\VarMStar\,$g shifts upward to tens of grams once taper and asymmetry
rescale the $\mathrm{Wo}$--mass mapping---whereas the hummingbird and neonatal
tests turn only on whether $\mathrm{Wo}_0$ crosses $\sqrt{3}$ and are therefore
robust to this refinement.} \textbf{The Hummingbird Test: A Falsifiable
Prediction.} A critical experimental test of the $\mathrm{Wo}$-driven transition
arises in high-performance avian physiology. Despite tiny body mass, extreme
heart rate can push small organisms into the wave-dominated regime.

\begin{table}[H]
\centering
\caption{Predictive verification of the Womersley-driven transition under cardiac frequency override. \chA{Throughout, $\mathrm{Wo}_0$ denotes the Womersley number at the aortic root (generation~0, the largest vessel), consistent with the per-generation hierarchy of the Supplemental Material; the mouse and hummingbird entries follow from the mouse baseline via Eq.~\eqref{eq:wo_scaling}.}}
\label{tab:hummingbird}
\begin{tabular}{ccccc}
\toprule
Species & $M$ & $f_H$ (bpm) & $\mathrm{Wo}_0$ & Predicted $\alpha^*$ \\
\midrule
Mouse               & 20\,g  & 600  & 2.5  & 2.60 (transition) \\
\textbf{Hummingbird} & \textbf{4\,g} & \textbf{1000} & \textbf{1.76} & \textbf{2.60 (transition)} \\
Human               & 70\,kg & \VarHeartRateHuman & \chA{17.9} & $\VarAlphaStar$ (minimax) \\
\bottomrule
\end{tabular}
\end{table}

\noindent{\small\textit{Falsification criterion:} Despite a tiny body mass
($M=4\,$g), the hummingbird's extreme heart rate yields $\mathrm{Wo}_0 \approx
1.76$, placing it precisely at the incommensurability edge ($\mathrm{Wo}_c =
\sqrt{3} \approx 1.73$). The theory predicts a transition attractor $\alpha
\approx 2.60$, \textbf{not} $\alpha \approx 3.0$ (Murray viscous) as would be
expected for such a microscopic animal under pure body-mass scaling. This
prediction is directly testable via morphometric analysis of hummingbird
vascular casts and provides a decisive test of whether Womersley number, rather
than body size alone, controls branching architecture.}

(For the hummingbird, scaling from the mouse baseline yields $\mathrm{Wo}_0 =
2.5 \times (4/20)^{3/8} \times (1000/600)^{1/2} \approx 1.76$ according to the
scaling of Eq.~\eqref{eq:wo_scaling}. \textbf{\chE{At present:}} To our
knowledge, no quantitative morphometric data for hummingbird coronary arteries
are currently available in the literature.)

\subsection{Independent Datasets}
To further demonstrate that the branching exponent $\alpha^*$ is not an isolated
feature of the Kassab porcine coronary dataset~\cite{kassab1993},
Table~\ref{tab:multivalid} summarizes validation measurements across independent
vascular and physiological networks.

\begin{table}[H]
\centering
\caption{\textbf{Multi-Organ and Multi-Kingdom Validation of the Minimax
Attractor.} Comparison of observed branching exponents across different
physiological systems and regimes. The minimax attractor $\alpha^* \approx
\VarAlphaStar$ characterizes pulsatile mammalian arterial beds, while the boundary
limits $\alpha_w \approx 2.0$ and $\alpha_t \approx \chA{\VarAlphaT}$ emerge in
developmental and static-transport systems, respectively.}
\label{tab:multivalid}
\vspace{0.5em}
\begin{tabular*}{\textwidth}{@{\extracolsep{\fill}} l c c c c @{}}
\toprule
System [Ref.] & Observed $\alpha$ & Regime & $d$ & Basis \\
\midrule
Cerebral Arteries \cite{rossitti1993}  & $2.50$--$2.90$ & Minimax ($\alpha^*$) & 3D & Pulsatile Isotropy \\
Coronary Arteries \cite{kassab1993}    & $2.60$--$2.80$ & Minimax ($\alpha^*$) & 3D & Pulsatile Isotropy \\
Human Retina (Healthy) \cite{luo2017}  & $2.0$--$2.6$   & Wave ($\alpha_w$)    & 2D & Planar Isotropy \\
Plant Xylem (Vines) \cite{mcculloh2003}& $2.90$--$3.00$ & Static ($\alpha_t$)  & 3D & Viscous Min. \\
\bottomrule
\end{tabular*}
\end{table}

The qualitative consistency of these results across fundamentally different
organs and kingdoms provides supporting evidence for the generality of the
incommensurability principle. An expanded validation across \chA{27} distinct
biological and synthetic networks is provided in the Supplemental Material
(Section S6), showing that the $\alpha^* \approx \VarAlphaStar$ attractor
appears consistently in pulsatile mammalian arterial systems.

\chRB{The breadth of this validation set is itself the central empirical test of
the principle, and not merely a catalogue of consistent cases. The
Incommensurability Principle does not predict a single universal exponent; it
predicts that the realized exponent must track \emph{which incommensurable cost
channels are physically active} in a given network. This yields a
\chRB{pre-specified}, cross-kingdom sorting rule that the 27 systems test
directly. (i)~Networks with no pulsatile wave channel---botanical xylem (lianas,
conifers, angiosperms), the bronchial and tracheal airways, leaf venation,
slime-mold and mycelial transport webs, and quasi-static venous return---should
occupy the pure viscous attractor $\alpha_t \to 3.0$; all ten such systems in
the catalogue fall within $2.80$--$3.15$. (ii)~Pulsatile mammalian arterial
beds, in which both channels are active, should occupy the minimax attractor
$\alpha^* \approx \VarAlphaStar$ (cerebral, coronary, pulmonary).
(iii)~Wave-dominated or planar networks---the $d=2$ retinal trees and the large
carotid and conduit segments---should approach the wave attractor $\alpha_w
\approx \VarAlphaW$. This sorting holds across three biological kingdoms
(animal, plant, and fungal/protist), both embedding dimensions, and pathological
states (hypertensive and tumour microvasculature shifting toward the wave-skewed
limit)---a span that no single-attractor theory, Murray's law ($\alpha=3$) or
the area-preserving law ($\alpha=2$) alone, can reproduce. Its discriminating
power is precisely that it is falsifiable in two directions: a non-pulsatile
network (plant, fungal, or airway) measured at $\alpha \approx \VarAlphaStar$,
or a high-Womersley pulsatile conduit measured at the viscous $\alpha \approx
3.0$, would each contradict the regime-sorting prediction. We regard the
systematic extension of this test---to taxa with divergent cardiac duty cycles
(avian and reptilian circulations), to invertebrate open circulatory systems,
and to engineered organ-on-chip networks with programmable pulsatility---as the
most informative direction for broader empirical validation, because each adds
an independent point at which the active-channel sorting could fail.}

\textbf{Quantitative validation with high-resolution datasets.} To move beyond
qualitative consistency, we perform direct quantitative comparison with all
available high-resolution morphometric datasets at the generation-by-generation
level demonstrated in Kassab (1993). Historically, three independent classical
morphometric studies have provided the necessary resolution:

\begin{enumerate}
    \item
\textbf{Kassab (1993)}: Porcine coronary, $\alpha_{\mathrm{obs}} = 2.60$--$2.80$
vs.\ $\alpha^*_{\mathrm{pred}} = \VarAlphaStarModel$ (symmetric minimax, this
work)
    \item
\chE{\textbf{Jiang, Kassab \& Fung (1994)}: Rat pulmonary ($M=300\,\mathrm{g}$),
$\alpha_{\mathrm{obs}} = 2.60$--$2.75$ vs.\ $\alpha^*_{\mathrm{pred}} = 2.68$}
    \item
\textbf{Huang (1996)}: Human pulmonary ($M=\VarMassHuman\,$g),
$\alpha_{\mathrm{obs}} = 2.60$--$2.85$ vs.\ $\alpha^*_{\mathrm{pred}} =
\VarAlphaStarModel$
\end{enumerate}

\chE{All three datasets exhibit quantitative agreement with theoretical
predictions. Collectively, the three datasets span about two orders of magnitude
in body mass ($300\,\mathrm{g}$ to $70\,\mathrm{kg}$) and two organ systems
(coronary, pulmonary). Detailed calculations are provided in the Supplemental
Material, Sections S6.3.3--S6.3.5.}

\textbf{Remaining limitations.} While quantitative validation has historically
been restricted to these three classical morphometric datasets, the majority of
\chA{the 27 systems in the extended validation set (Supplemental Material,
Section S6)} exhibit wide observed ranges that overlap with multiple theoretical
attractors. This data limitation is rapidly being resolved by next-generation
whole-organ imaging databases: for instance, complete 3D microCT reconstructions
of the mouse cerebral angiome~\cite{quintana2019}, anatomically detailed
finite-element pulmonary models~\cite{burrowes2005}, and whole-organ isotropic
synchrotron datasets (using the original Hierarchical Phase-Contrast Tomography,
HiP-CT, framework~\cite{walsh2021} and its whole-organ
applications~\cite{chourrout2026}) are providing high-resolution,
generation-by-generation parent-daughter connectivity maps. These advanced
datasets provide powerful, independent tests of our minimax scaling predictions.

\subsection{Dimensional Sensitivity}
\label{sec:dimensional_sensitivity}

A fundamental consequence of the minimax framework is that the dual thresholds
$\mathrm{Wo}_c^{\mathrm{fluid}}$ and $\mathrm{Wo}_c^{\mathrm{wave}}$ are
sensitive to the embedding dimensionality $d$. For a network constrained to a
two-dimensional manifold (such as the retinal vasculature or leaf venation), the
requirement of isotropic power flow simplifies to planar rotational invariance
$SO(2)$. Physically, this "Symmetry-Preservation Hypothesis" identifies the
ground state of the endothelial mechanotransduction machinery: isotropic shear
stress minimizes local stress gradients on the cell membrane, providing a
mechanical null-point for homeostatic stability. Under these conditions: the
emergent phase angle of the complex fluid impedance must satisfy $\tan \phi =
d-1 = 1$, yielding an isotropic phase of $\phi_{\mathrm{iso}} = 45^\circ$.
Substituting this into the fluid admittance expansion $\tan\phi \approx
6/\mathrm{Wo}^2$ yields a higher critical fluid threshold:
$\mathrm{Wo}_c^{\mathrm{fluid}}(d=2) = \sqrt{6} \approx \VarWoCTwo$ (compared to
$\sqrt{3} \approx \VarWoC$ in 3D; see Supplemental Material, Section S1 for the
numerical verification).

Since $M^* \propto (\mathrm{Wo}_c^{\mathrm{fluid}})^4$, the critical fluid
transition mass for 2D networks is shifted upward to $M^*_{d=2} \approx
\VarMStarTwo\,$g, approximately four times the 3D value. However, as
demonstrated in the Phase Decoupling derivation, the corresponding wave
threshold $\mathrm{Wo}_c^{\mathrm{wave}}$ collapses identically to $0$ in 2D.
This striking mathematical result implies that the wave is \emph{always} above
its threshold in planar geometries, \chA{placing} planar networks (retinal
vasculature, embryonic membranes) into the wave-dominated regime regardless of
their small physical mass or local Womersley number.

This shift is directly supported by recent 3D whole-embryo mappings
\cite{forjaz2026}, which reveal a sharp divergence between 3D space-filling
vascular transport ($\alpha \approx 3.0$) and 2D planar neural structures
($\alpha \approx 2.0$). \chA{Additional support} is provided by the human retina
\cite{luo2017}, where retinal arteries range from $\alpha \approx 2.1$ (small
vessels) to $2.7$ (large vessels), consistent with the 2D wave-dominated regime
prediction.

Interestingly, while retinal diameters follow the wave-attractor ($\alpha
\approx \VarAlphaWTwo$), the mean branching angles in the RBAD dataset
($\VarAngleRetinalObs^\circ \pm \VarAngleRetinalErr^\circ$ (SE, $n{=}342$
bifurcations; SD $\VarAngleRetinalSD^\circ$)) remain closer to the Murray-limit
($\alpha \approx \VarAlphaT$). This discrepancy, termed the "Angle-Diameter
Paradox," reveals a fundamental Mechanical-Fluidic Decoupling: diameters respond
to wave-reflection constraints (2D), while branching angles are anchored by the
mechanical equilibrium of wall tensions (3D). While the global network
architecture (diameters) is \chA{biased toward the wave regime} to minimize
reflection-induced dissipation, the local junctional geometry (angles) is
governed by the mechanical equilibrium of wall tensions. As long as the vascular
wall operates in the Lam\'{e}-static regime, the branching angles will remain
anchored to the Murray-limit, even as the fluidic impedance collapses toward the
area-preserving limit. This "structural tension" is \chA{one direct empirical}
\chA{signature of the Scaling Conflict Bound}: the incommensurability between
structural integrity and fluidic efficiency prevents a single, unified scaling
exponent from governing all morphometric degrees of freedom under symmetric
local optimization.

In mammalian arterial beds (Coronary, Cerebral, Renal), the convergence towards
the minimax reflects the balance between metabolic maintenance and
wave-reflection integrity. In contrast, plant xylem networks, which lack a
pulsatile pump ($\eta \to 0$), converge towards the static Murray-limit ($\alpha
\approx \VarAlphaT$). Conversely, developmental networks and veins, where
metabolic cost is secondary to signal propagation, gravitate towards the wave
limit ($\alpha \approx \VarAlphaWTwo$). This mapping \chA{supports the
interpretation} that biological transport architecture is governed by a
universal \emph{minimax attractor} that shifts predictably across physical
regimes.
